\scanpage{71}
a) $\exists n\in\mathbb N$ tel que $T_n\ne\varnothing$

\inserted{b) $\forall m\in\mathbb N$ et $v\in T_m$, on a $v(E)\in C_m$}

$\unsure{c|b})$ Si \deleted{$\unsure{u\in T_n}$} $(u:E\to R^n)\in T_n$, et si $v:E\lhook\joinrel\longrightarrow R^m$, alors $v\in T_m$ ssi $1^\circ)$ $v(E)\in C_m$ et $2^\circ)$
\[
vu^{-1}:\underset{\cap\,R^n}{u(E)}
\longrightarrow\underset{\cap\,R^m}{v(E)}
\]
est un $C_*$-morphisme

\margin{Cette structure peut aussi être définie par l'ens $\Sigma$ des $C_*$-morphismes de $E$ dans $R$, avec l'axiome suivant
a) $\exists J\subset\Sigma$, partie finie, telle que $E\to R^J$ can. soit injectif, et ait une image $X\in C_J$.
b) On peut choisir $J$ tel que $\forall f:E\to R$, on ait $f\in\Sigma$ ssi le composé
\diagram{\begin{tikzcd}[ampersand replacement=\&]
X\arrow[r,"{\sim}"]\&E\arrow[r,"{f}"]\&R
\end{tikzcd}}
est $C_*$-morphisme,
\inserted{inverse de $E\overset{\sim}{\to}X$}.}

\uline{Ex} Si $n\in\mathbb N$, $X\subset R^n$ avec $X\in C_n$, alors $X$ est muni d'une $C_*$-structure canonique.

\uline{Déf} Soient $X,Y$ deux $C_*$-modèles. Une application $f:X\to Y$ est dite un $C_*$-morphisme si ---

\uline{Prop} Les \deleted{dé} $C_*$-modèles, et $C_*$-morphismes, forment une catégorie (pour la composition ensembliste des morphismes).
Dans cette catégorie, les produits \deleted{finis} de deux objets existent.
Un morphisme de $C_*$-modèles est un iso. ssi il est bijectif.

\uline{Prop} Conditions équivalentes sur $C_*$

a) \deleted{$\unsure{C_0\ne\varnothing}$} $\exists n\in\mathbb N$ et $X\in C_n$, tel que $X\ne\varnothing$

b) $R^0\in C_0$

c) La catégorie des $C_*$-modèles admet des objets de cardinal $1$ (i.e. sur un ens. à $1$ élément, existe une structure de $C_*$-modèle)

Alors les \deleted{\unsure{modè}} $C_*$-modèles \unsure{uniponctuels} sont des objets finaux de la catégorie des $C_*$-modèles.

Dorénavant, on supposera cette condition, sous la forme plus forte

\textcircled{5} \deleted{$R^0\in C_0$}
\inserted{$\forall n\in\mathbb N$, on a}
\[
R^n=\bigcup_{X\in C_n}X
\]

\margin{On définit la notion de $C_*$-partie d'un $C_*$-modèle, de façon unique, de façon que ça commute au ``transport de structure'' et ait Si $X\in C_n$ ($n\in\mathbb N$), alors pour $Y\subset X$, $Y$ est une $C_*$-partie du $C_*$-modèle $X$ ssi $Y\in C_n$.
Ceci dit, pour un morphisme de modèles, l'image directe ou inverse d'une $C_*$-partie est itou.
Notion de $C_*$-str. induite.}

\margin{72}
